\section{Moving embeddings and the edge phase space}\label{sec:moving-boundary}

\subsection{Pulling back the complete action}\label{subsec:pullback-action}

A fixed coordinate wall is not a gauge-invariant notion. Introduce an embedding $X:M_0\to M$ of a reference region and the field-space displacement

\begin{align}
\chi^\mu:=\delta X^\mu\circ X^{-1}.
\end{align}

For every spacetime tensor, the covariant field-space variation is

\begin{align}
\Delta_Xg:=\delta g+\mathcal L_\chi g,
\qquad
\delta(X^*g)=X^*(\Delta_Xg). \label{eq:delta-x}
\end{align}

We define the extended theory by pulling back the complete finite action,

\begin{align}
S_{{\rm p},X}[g,X]
:=S_{\rm p}[X^*g;M_0],
\end{align}

and similarly define its potential by

\begin{align}
\Theta_{{\rm p},X}[g,X;\delta]
:=\Theta_{\rm p}[X^*g;\delta(X^*g)]. \label{eq:extended-potential}
\end{align}

Equation (\ref{eq:extended-potential}) includes the shape variation of the GHY surfaces, normals, counterterms, Harlow--Wu endpoint terms, and Hayward joints. It is not a metric-only potential with an embedding term appended afterward.

Let $\widehat w$ act by

\begin{align}
I_{\widehat w}\delta g=\mathcal L_wg,
\qquad
I_{\widehat w}\chi=-w.
\end{align}

Then $I_{\widehat w}\Delta_Xg=0$. Hence

\begin{align}
(\mathcal L_wg,-w\circ X)
\in\ker\Omega_{{\rm p},X}. \label{eq:extended-degeneracy}
\end{align}

This is the precise extended-space statement of gauge invariance. A metric-only transformation at fixed $X$ and an embedding-only transformation at fixed $g$ carry opposite surface charges; only their diagonal combination is a degeneracy \cite{Donnelly:2016auv,Speranza:2017guc}.

\subsection{Wall embedding momentum}\label{subsec:embedding-momentum}

On a moving timelike wall, the source one-form becomes

\begin{align}
\mathcal B_X=-\frac12\Pi^{ij}\Delta_X\gamma_{ij}.
\end{align}

Decompose $\chi=\chi_\perp n+\chi^ie_i$. Variation of the boundary term gives, up to endpoint total derivatives,

\begin{align}
P_j&=D_i\Pi^i{}_j,\\
P_\perp&=-\Pi^{ij}K_{ij}.
\end{align}

These are the Brown--York contributions, not yet the complete embedding momentum. Direct variation of the bulk domain supplies the constraints, giving

\begin{align}
P_j^{\rm tot}
&=D_i\Pi^i{}_j
-\frac{\sqrt{-\gamma}}{\kappa_{\rm p}^{2}}
n_\mu\mathcal E^{\mu\nu}e_{j\nu}, \label{eq:tangential-momentum}\\
P_\perp^{\rm tot}
&=-\Pi^{ij}K_{ij}
-\frac{\sqrt{-\gamma}}{\kappa_{\rm p}^{2}}
n_\mu n_\nu\mathcal E^{\mu\nu}. \label{eq:normal-momentum}
\end{align}

The cosmological Einstein tensor $\mathcal E^{\mu\nu}$ therefore completes the Brown--York work off shell. On shell, (\ref{eq:tangential-momentum}) and (\ref{eq:normal-momentum}) reduce to the familiar boundary momentum and normal work.

For a non-orthogonal moving joint, varying the Hayward term gives the edge two-form

\begin{align}
\Omega_J
=\frac{\sigma_J}{\kappa_{\rm p}^{2}}
\int_J\Delta_X\eta\wedge\Delta_X\sqrt q . \label{eq:hayward-pair}
\end{align}

Thus the canonical joint pair is area density and finite relative boost angle, not area and surface gravity. Equation (\ref{eq:hayward-pair}) is the action-derived realization of the area--boost edge pair \cite{Takayanagi:2019tvn}.

\subsection{The horizon cocycle as a change of section}\label{subsec:section-change}

Let $s_0$ be the fixed-embedding section. A metric representative change $h\mapsto h+\mathcal L_vG$ at fixed embedding is not the full gauge transformation. In the extended space,

\begin{align}
(h,0)\sim(h+\mathcal L_vG,-v). \label{eq:section-change}
\end{align}

By (\ref{eq:extended-degeneracy}), the change of metric canonical energy between the two fixed-embedding descriptions is cancelled by the embedding and joint contribution. Therefore

\begin{align}
\int_\gamma\Upsilon_{\rm p}[h,v]
=E_{{\rm can,p}}[h+\mathcal L_vG]
-E_{{\rm can,p}}[h]
\end{align}

is naturally the cocycle associated with changing the section of the extended phase space. This interpretation derives the required compensation from the action without asserting that $\Upsilon_{\rm p}$ is a globally integrable scalar functional on a nonlinear edge space.

One can always cancel a closed finite-dimensional two-form $F$ by adjoining an auxiliary cotangent graph. For constant coefficients,

\begin{align}
\Theta_{\rm aux}=-\frac12F_{IJ}a^I\delta a^J,
\qquad
\delta\Theta_{\rm aux}=-F.
\end{align}

This local symplectic fact is not used as a physical input in our theorem. In particular, adding a scalar boundary functional $w$ cannot change $\Omega$, since $\delta^2w=0$. The actual cancellation in (\ref{eq:section-change}) comes from the pulled-back gravitational action and its embedding variation.

\subsection{Nonuniform compact polarization}\label{subsec:compact-polarization}

The moving construction is nonempty beyond a rigid wall. In an exact BTZ collar consider

\begin{align}
\mathrm ds^2
=\mathrm d\rho^2-s_\xi^2\sinh^2\rho\,\mathrm d\lambda^2
+R(\phi)^2\cosh^2\rho\,\mathrm d\phi^2,
\end{align}

with a stationary inner wall $\rho=E(\phi)>0$. Define

\begin{align}
\vartheta
=\arctan\left(\frac{E'}{R\cosh E}\right).
\end{align}

The complete Brown--York source is exactly

\[\begin{aligned}
\mathcal B_X
={}&-\delta\left\{
\frac{s_\xi}{\kappa_{\rm p}^{2}}
\left[R\cosh^2E+E'\cosh E\,\vartheta\right]
\right\}\\
&+\partial_\phi\left(
\frac{s_\xi\cosh E}{\kappa_{\rm p}^{2}}
\vartheta\,\delta E
\right). \label{eq:nonuniform-polarization}
\end{aligned}\]

The last term integrates to zero on the circle, so the stationary smooth $(R,E)$ family is a Lagrangian moving-wall polarization with arbitrary nonuniform normal displacement. This is a genuine infinite-dimensional compact family. It does not prove a simultaneous time-dependent wavy-wall and tilted-joint polarization, which lies outside the present theorem.
